\chapter*{参考文献与延伸阅读}\addcontentsline{toc}{chapter}{参考文献与延伸阅读}
\begin{enumerate}[leftmargin=*,itemsep=2pt,parsep=0pt]\small
\item Bintu, L. et al. (2005). Transcriptional regulation by the numbers: models. \emph{Current Opinion in Genetics \& Development}, 15, 116--124. DOI: \href{https://doi.org/10.1016/j.gde.2005.02.007}{\nolinkurl{10.1016/j.gde.2005.02.007}}.
\item Gardner, T. S., Cantor, C. R., \& Collins, J. J. (2000). Construction of a genetic toggle switch in Escherichia coli. \emph{Nature}, 403, 339--342. DOI: \href{https://doi.org/10.1038/35002131}{\nolinkurl{10.1038/35002131}}.
\item Elowitz, M. B. \& Leibler, S. (2000). A synthetic oscillatory network of transcriptional regulators. \emph{Nature}, 403, 335--338. DOI: \href{https://doi.org/10.1038/35002125}{\nolinkurl{10.1038/35002125}}.
\item Shen-Orr, S. S. et al. (2002). Network motifs in the transcriptional regulation network of Escherichia coli. \emph{Nature Genetics}, 31, 64--68. DOI: \href{https://doi.org/10.1038/ng881}{\nolinkurl{10.1038/ng881}}.
\item Alon, U. (2007). Network motifs: theory and experimental approaches. \emph{Nature Reviews Genetics}, 8, 450--461. DOI: \href{https://doi.org/10.1038/nrg2102}{\nolinkurl{10.1038/nrg2102}}.
\item Kuznetsov, Y. A. Andronov--Hopf bifurcation. \emph{Scholarpedia}. \href{https://www.scholarpedia.org/article/Andronov-Hopf_bifurcation}{Hopf 定理与第一 Lyapunov 系数公式}。
\item Utrecht University. The Poincar\'e--Bendixson Theorem. \href{https://webspace.science.uu.nl/~hanss102/inlds224/PoincareBendixsonTheorem.pdf}{平面极限集定理与证明讲义}。
\end{enumerate}
